% !TEX root = main.tex
\section{ParComm: Collection Commitment and Opening}
\label{sec:parcomm}

ParComm commits to a collection of matrices through one BVA.  At opening time,
it compiles the independent left and right evaluation points into one TBV,
expresses all matrix evaluations as a single BVA--TBV inner product, and applies
Packed EO-LiteBullet to reduce this inner product to an MLE claim on the
committed BVA.  The compilation is proved from one commitment to the original
points and the registered shape-only matrices: lookup proofs bind the local
points and their source rows, while one prefix recurrence derives the TBV
scales directly from the selected point coordinates.
The remaining TBV evaluation is checked by one prefix recurrence.

\subsection{Recursive Matrix Splitting}
\label{sec:matrix-splitting}

For each $\tau\in[t]$, let
$\bm C_\tau\in\Field^{\mu_\tau\times\nu_\tau}$, and let
\[
  \bm x_{l,\tau}\in\Field^{\lceil\log_2\mu_\tau\rceil},
  \qquad
  \bm x_{r,\tau}\in\Field^{\lceil\log_2\nu_\tau\rceil}
\]
be its left and right evaluation points.  We write
\[
  \bm l_\tau:=\tensorspan{\bm x_{l,\tau}},
  \qquad
  \bm r_\tau:=\tensorspan{\bm x_{r,\tau}}.
\]
Let $\overline{\bm C}_\tau$ denote the zero extension of $\bm C_\tau$ to the
dimensions of $\bm l_\tau$ and $\bm r_\tau$.  The claimed evaluation is
\begin{equation}
\label{eq:individual-evaluations}
\begin{aligned}
  u_\tau
  &=\bm l_\tau\overline{\bm C}_\tau\bm r_\tau^{\top}\\
  &=\left\langle
      \vectorize{\overline{\bm C}_\tau},
      \tensorspan{\bm x_{r,\tau}\|\bm x_{l,\tau}}
    \right\rangle .
\end{aligned}
\end{equation}
The order $\bm x_{r,\tau}\|\bm x_{l,\tau}$ follows from row-major
vectorization.

\begin{lemma}[Recursive dyadic splitting]
\label{lemma:mat-splitting}
A $\mu\times\nu$ matrix can be partitioned into at most
\[
  (\lfloor\log_2\mu\rfloor+1)
  (\lfloor\log_2\nu\rfloor+1)
\]
  aligned submatrices whose two dimensions are powers of two.  Hence, $t$
matrices of total size $N$ produce
$T=O(t\log^2 N)$ such submatrices.
\end{lemma}

\begin{proof}
Let $\alpha=\lfloor\log_2\mu\rfloor$ and
$\beta=\lfloor\log_2\nu\rfloor$.  First separate the top $2^\alpha$ rows.
Splitting their columns according to the binary expansion of $\nu$ gives at
most $\beta+1$ power-of-two submatrices.  Next, in the remaining lower part,
separate the leftmost $2^\beta$ columns.  Splitting their rows according to the
binary expansion of $\mu-2^\alpha$ gives at most $\alpha$ submatrices.

It remains to process the lower-right block.  Both of its logarithmic
dimensions are at most $\alpha-1$ and $\beta-1$, respectively.  Recursing on
this block therefore gives, by induction, at most $\alpha\beta$ further
submatrices.  The total is at most
\[
  (\beta+1)+\alpha+\alpha\beta
  =(\alpha+1)(\beta+1).
\]
The recursion stops when one dimension is one, in which case the other
dimension is partitioned by its binary expansion.  Every resulting interval is
aligned to its own length because it starts after a sum of larger powers of
two.  Summing the bound over all matrices proves the second claim.
\end{proof}

Apply this deterministic splitting rule to all matrices.  Let
$\bm a^{(j)}$ be the row-major vectorization of the resulting submatrices,
ordered by non-increasing size, with ties broken by the original matrix,
row, and column positions.  Define
\[
  k_j:=|\bm a^{(j)}|,
  \qquad j\in[T].
\]
Every $k_j$ is a power of two and
\[
    \sum_{j=1}^{T}k_j=N.
\]
Let
\[
  K:=\mathsf{next\_power\_of\_two}(N),
  \qquad
  m:=\log_2K,
\]
and form the BVA
\begin{equation}
\label{eq:matrix-bva}
  \bm a
  :=\bm a^{(1)}\|\cdots\|\bm a^{(T)}\|\bm0^{K-N}.
\end{equation}
Because the bucket sizes are non-increasing powers of two, each bucket is
aligned to its own size.  The public layout is
$\mathcal L=(K,N,T,\{k_j\}_{j\in[T]})$, and the prover commits once:
\begin{equation}
\label{eq:dory-bva-commit}
  C_a=\pc.\Com(\pp,\bm a).
\end{equation}
Only the $N$ matrix entries are stored; the final zero suffix is implicit.


\subsection{Compiling Matrix Points into a TBV}
\label{sec:point-compilation}

For clarity, let the BVA cover size be a power of two and write
$m:=\log_2K$.  Consider matrices
$\{\bm C_\tau\}_{\tau\in[t]}$ with evaluation claims
\[
    u_\tau
    =
    \bm l_\tau^{\top}\bm C_\tau\bm r_\tau,
    \qquad
    \bm l_\tau=\tensorspan{\bm x_{l,\tau}},
    \qquad
    \bm r_\tau=\tensorspan{\bm x_{r,\tau}}.
\]
Lemma~\ref{lemma:mat-splitting} partitions these matrices into
\[
    T\le t(m+1)^2
\]
aligned submatrices whose row and column dimensions are powers of two.
We compile the evaluation points of these submatrices into the scales
$\eta_j$ and local points $\bm z_j$ of a single TBV.

\mysmallskip\noindent
\textbf{Original evaluation points.}
Pad every point to length $m$ and arrange the points row-wise:
\begin{equation}
\label{eq:point-matrix}
\bm X_l=
\begin{pmatrix}
\bm x_{l,1}^{\top}\|\bm0^{\top}\\[-1mm]
\vdots\\[-1mm]
\bm x_{l,t}^{\top}\|\bm0^{\top}
\end{pmatrix},
\qquad
\bm X_r=
\begin{pmatrix}
\bm x_{r,1}^{\top}\|\bm0^{\top}\\[-1mm]
\vdots\\[-1mm]
\bm x_{r,t}^{\top}\|\bm0^{\top}
\end{pmatrix}.
\end{equation}
Thus,
\[
    \bm X_l,\bm X_r\in\Field^{t\times m}.
\]
We prepend one fixed zero column and define
\begin{equation}
\label{eq:combined-point-matrix}
    \bm X
    :=
    \concat{\bm0_{t\times1},\bm X_r,\bm X_l}
    \in\Field^{t\times(2m+1)}.
\end{equation}
Column $0$ is the fixed zero entry, columns $1,\ldots,m$ contain the
right points, and columns $m+1,\ldots,2m$ contain the left points.  The order
$\bm X_r\|\bm X_l$ follows row-major vectorization.  The prover commits to
the point matrix as
\[
    C_X=\pc.\Com(\pp,\bm X).
\]

\mysmallskip\noindent
\textbf{Shape metadata.}
Index the $j$-th split submatrix by
\[
    (\tau_j,\alpha_{l,j},\gamma_{l,j},
             \alpha_{r,j},\gamma_{r,j}),
\]
where $\tau_j$ identifies its source matrix,
$(\alpha_{l,j},\alpha_{r,j})$ is its row and column offset, and its shape is
\[
    2^{\gamma_{l,j}}\times2^{\gamma_{r,j}}.
\]
Let $\bm\delta_{l,j}$ and $\bm\delta_{r,j}$ be the $m$-bit little-endian
decompositions of $\alpha_{l,j}$ and $\alpha_{r,j}$, respectively.
Alignment gives
\begin{equation}
\label{eq:split-alignment}
\begin{aligned}
    \delta_{l,j,h}&=0,
        &&h\in[\gamma_{l,j}],\\
    \delta_{r,j,h}&=0,
        &&h\in[\gamma_{r,j}].
\end{aligned}
\end{equation}

Arrange the metadata as
\[
\begin{aligned}
    \bm\tau&=(\tau_1,\ldots,\tau_T)^\top,\\
    \bm\alpha_l&=(\alpha_{l,1},\ldots,\alpha_{l,T})^\top,
    &\bm\alpha_r&=(\alpha_{r,1},\ldots,\alpha_{r,T})^\top,\\
    \bm\gamma_l&=(\gamma_{l,1},\ldots,\gamma_{l,T})^\top,
    &\bm\gamma_r&=(\gamma_{r,1},\ldots,\gamma_{r,T})^\top.
\end{aligned}
\]
The primary shape metadata is
\begin{equation}
\label{eq:shape-metadata}
    \bm\alpha
    :=
    \concat{\bm\tau,\bm\alpha_l,\bm\alpha_r},
    \qquad
    \bm\gamma
    :=
    \concat{\bm\gamma_l,\bm\gamma_r}.
\end{equation}
Thus, $\bm\alpha\in\Field^{T\times3}$ and
$\bm\gamma\in\Field^{T\times2}$.  In particular, the source-matrix
assignment $\bm\tau$ is included in $\bm\alpha$ and requires no separate
persistent commitment.

The public size $k_j$ of the corresponding BVA bucket satisfies
\begin{equation}
\label{eq:bucket-size-gamma}
    \gamma_{l,j}+\gamma_{r,j}
    =
    \log_2 k_j.
\end{equation}

\mysmallskip\noindent
\textbf{Derived matrices.}
The local evaluation point of the $j$-th submatrix is
\begin{equation}
\label{eq:matZfromX}
    \bm z_j
    =
    \bm X_{r,\tau_j,[1:\gamma_{r,j}]}^\top
    \|
    \bm X_{l,\tau_j,[1:\gamma_{l,j}]}^\top.
\end{equation}
It contains exactly the coordinates that vary within the submatrix and has
length
\[
    |\bm z_j|
    =
    \gamma_{r,j}+\gamma_{l,j}
    =
    \log_2k_j.
\]

Let $\bm\sigma_j\in\Field^{2m}$ record the corresponding column positions
in $\bm X$:
\begin{equation}
\label{eq:sigma-row}
\bm\sigma_j
=
\bigl(
1,\ldots,\gamma_{r,j},
m+1,\ldots,m+\gamma_{l,j},
0,\ldots,0
\bigr).
\end{equation}
The zero entries select the fixed zero column of $\bm X$.

Arrange the derived vectors row-wise:
\begin{equation}
\label{eq:compiled-point-matrices}
\bm\Sigma=
\begin{pmatrix}
\bm\sigma_1^\top\\[-1mm]
\vdots\\[-1mm]
\bm\sigma_T^\top
\end{pmatrix},
\quad
\bm D=
\begin{pmatrix}
(\bm\delta_{r,1}\|\bm\delta_{l,1})^\top\\[-1mm]
\vdots\\[-1mm]
(\bm\delta_{r,T}\|\bm\delta_{l,T})^\top
\end{pmatrix},
\quad
\bm Z=
\begin{pmatrix}
\bm z_1^\top\|\bm0^\top\\[-1mm]
\vdots\\[-1mm]
\bm z_T^\top\|\bm0^\top
\end{pmatrix}.
\end{equation}
All three matrices have dimensions $T\times2m$.  Only the first
$\gamma_{r,j}+\gamma_{l,j}$ entries of row $j$ of $\bm Z$ form the local
point $\bm z_j$; the remaining entries are zero padding.

Finally, copy the source point row associated with each bucket:
\begin{equation}
\label{eq:selected-point-matrix}
\bm X'
=
\begin{pmatrix}
(\bm X_r\|\bm X_l)_{\tau_1,*}\\[-1mm]
\vdots\\[-1mm]
(\bm X_r\|\bm X_l)_{\tau_T,*}
\end{pmatrix}
\in\Field^{T\times2m}.
\end{equation}
The matrices $\bm\Sigma$ and $\bm D$ depend only on the deterministic split,
not on an evaluation point.  They are therefore committed and proved once
during registration together with $\bm\alpha$ and $\bm\gamma$.  Every opening
reuses the same commitments $C_\Sigma$ and $C_D$.  Their consistency with
$\bm\gamma$ and $\bm\alpha$ is established by the following lemma.

\begin{lemma}[Consistency of the Shape-Derived Matrices]
\label{lemma:shape-derived-consistency}
Suppose that
\[
    0\le\gamma_{l,j},\gamma_{r,j}\le m,
    \qquad
    0\le\alpha_{l,j},\alpha_{r,j}<N
\]
for every $j\in[T]$.  Let
$\bm\Sigma$, $\bm D$, $\bm\alpha$, and $\bm\gamma$ be fixed before
$\vartheta\ChooseRandom\Field\setminus\{0,1\}$ is sampled.

Consistency between $\bm\Sigma$ and $\bm\gamma$ is checked simultaneously
for all rows by the following matrix equation:
\begin{equation}
\label{eq:Sigma-gamma-polynomial}
\begin{aligned}
&(1-\vartheta)
\Bigg[
(1-\vartheta)\bm\Sigma
(1,\vartheta,\ldots,\vartheta^{2m-1})^\top\\
&\qquad
+(\bm\gamma_r-m\bm1_T)
    \circ\vartheta^{\bm\gamma_r}\\
&\qquad
+(m\bm1_T+\bm\gamma_l)
    \circ\vartheta^{\bm\gamma_r+\bm\gamma_l}
\Bigg]\\
&\hspace{12mm}
=
\bm1_T-\vartheta^{\bm\gamma_r+\bm\gamma_l},
\end{aligned}
\end{equation}
where exponentiation of a vector is defined componentwise; for example,
\[
\vartheta^{\bm\gamma_r}
=
\bigl(
\vartheta^{\gamma_{r,1}},
\ldots,
\vartheta^{\gamma_{r,T}}
\bigr)^\top.
\]
Except with probability at most
\[
    \frac{2m-1}{|\Field|-2},
\]
Eq.~\eqref{eq:Sigma-gamma-polynomial} simultaneously enforces, for every
$j\in[T]$,
\[
\bm\Sigma_{j,*}
=
\bigl(
1,\ldots,\gamma_{r,j},
m+1,\ldots,m+\gamma_{l,j},
0,\ldots,0
\bigr).
\]

Consistency between $\bm D$ and the offset columns of $\bm\alpha$ is
enforced by
\begin{equation}
\label{eq:D-alpha-consistency}
\begin{aligned}
    \bm D\circ(\bm1-\bm D)
    &=\bm0,\\
    \bm D_{*,[1:m]}
    (1,2,\ldots,2^{m-1})^\top
    &=\bm\alpha_r,\\
    \bm D_{*,[m+1:2m]}
    (1,2,\ldots,2^{m-1})^\top
    &=\bm\alpha_l.
\end{aligned}
\end{equation}
These equations uniquely determine $\bm D$ as the little-endian bit
decomposition of $(\bm\alpha_r,\bm\alpha_l)$.
\end{lemma}

\begin{proof}
For the row prescribed by Eq.~\eqref{eq:sigma-row}, we have
\[
\begin{aligned}
&
\left\langle
\bm\sigma_j,
(1,\vartheta,\ldots,\vartheta^{2m-1})
\right\rangle\\
&\quad=
\sum_{h=1}^{\gamma_{r,j}}
    h\vartheta^{h-1}
+
\sum_{h=1}^{\gamma_{l,j}}
    (m+h)\vartheta^{\gamma_{r,j}+h-1}.
\end{aligned}
\]
Substituting this expression into
Eq.~\eqref{eq:Sigma-gamma-polynomial} and telescoping consecutive
coefficients gives
\[
    1-\vartheta^{\gamma_{r,j}+\gamma_{l,j}},
\]
which proves completeness.

For soundness, let $\bm\sigma_j^\star$ denote the prescribed row and
subtract its valid equation from
Eq.~\eqref{eq:Sigma-gamma-polynomial}.  If
$\bm\sigma_j\neq\bm\sigma_j^\star$, then
\[
(1-\vartheta)^2
\left\langle
\bm\sigma_j-\bm\sigma_j^\star,
(1,\vartheta,\ldots,\vartheta^{2m-1})
\right\rangle
=0
\]
is a nonzero polynomial equation of degree at most $2m-1$ after excluding
the fixed root $\vartheta=1$.  The claimed bound follows from the
Schwartz--Zippel lemma.

Since $\bm\gamma$ is committed, the powers in
Eq.~\eqref{eq:Sigma-gamma-polynomial} are bound to their exponents by
lookup from the public table
\[
    \{(h,\vartheta^h):h=0,\ldots,2m\}.
\]
Equation~\eqref{eq:bucket-size-gamma} additionally binds
$\gamma_{r,j}+\gamma_{l,j}$ to the public bucket size.

For $\bm D$, the first equation of
Eq.~\eqref{eq:D-alpha-consistency} makes every entry Boolean.  The
remaining two equations reconstruct the numerical offsets from their
little-endian bits.  Since each offset lies in $[0,N)$, its $m$-bit
representation is unique.
\end{proof}

It remains to enforce alignment.  Since the nonzero entries of
$\bm\sigma_j$ are precisely
\[
    1,\ldots,\gamma_{r,j},
    \qquad
    m+1,\ldots,m+\gamma_{l,j},
\]
the indexed-zero constraints
\begin{equation}
\label{eq:alignment-indexed-zero}
    D_{j,\Sigma_{j,h}}=0
    \qquad
    \text{whenever }\Sigma_{j,h}\neq0
\end{equation}
are equivalent to Eq.~\eqref{eq:split-alignment}.  These constraints are
proved using the same indexed-read argument as the point lookup below.

\mysmallskip\noindent
\textbf{Compiling the local points.}
After the opening-time matrices $\bm Z$ and $\bm X'$ are fixed, the prover
uses the registered $\bm\Sigma$ in the following row-wise lookup:
\begin{equation}
\label{eq:coordinate-lookup}
\begin{aligned}
&
\left\{
Z_{j,h}+\vartheta\Sigma_{j,h}
:
h\in[2m]
\right\}\\
&\quad\subseteq
\{0\}
\cup
\left\{
X'_{j,h}+\vartheta h
:
h\in[2m]
\right\}.
\end{aligned}
\end{equation}
The challenge $\vartheta$ binds each value to its prescribed column.
A padded pair $(Z_{j,h},\Sigma_{j,h})=(0,0)$ reads the fixed zero entry.
Together with the cached registration check in
Eq.~\eqref{eq:Sigma-gamma-polynomial}, this proves that row $j$ of
$\bm Z$ contains exactly the local point prescribed by
$(\gamma_{r,j},\gamma_{l,j})$.

\mysmallskip\noindent
\textbf{Binding the source matrix.}
After receiving the batching challenge $\omega$, define
\[
\bm n=
\begin{pmatrix}
1\\[-1mm]
\vdots\\[-1mm]
t
\end{pmatrix},
\qquad
\bm\omega=
\begin{pmatrix}
1\\[-1mm]
\omega\\[-1mm]
\vdots\\[-1mm]
\omega^{t-1}
\end{pmatrix},
\qquad
\bm\omega'=
\begin{pmatrix}
\omega^{\tau_1-1}\\[-1mm]
\vdots\\[-1mm]
\omega^{\tau_T-1}
\end{pmatrix}.
\]
Once $\bm X'$ and $\bm\omega'$ are fixed, the verifier samples
$\theta\in\Field^*$ and sets
\[
    \bm\theta
    =
    (\theta^2,\theta^3,\ldots,\theta^{2m+1})^\top.
\]
The prover proves the multiset lookup
\begin{equation}
\label{eq:source-row-lookup}
\begin{aligned}
\left\{
\bm\tau+\theta\bm\omega'
+\bm X'\bm\theta
\right\}
\subseteq
\left\{
\bm n+\theta\bm\omega
+\concat{\bm X_r,\bm X_l}\bm\theta
\right\}.
\end{aligned}
\end{equation}
Here the braces denote the multisets of vector entries.
The random linear compression simultaneously binds, for every $j\in[T]$,
\[
    \bm X'_{j,*}
    =
    \concat{\bm X_r,\bm X_l}_{\tau_j,*},
    \qquad
    \omega'_j
    =
    \omega^{\tau_j-1}.
\]
Since $\bm\tau$ is the first column of the committed matrix $\bm\alpha$,
the source-matrix assignment is bound to the persistent shape encoding.

Equations~\eqref{eq:coordinate-lookup} and
\eqref{eq:source-row-lookup} can be summarized as
\begin{equation}
\label{eq:point-lookup}
\begin{aligned}
    (\bm X',\bm\Sigma)
    &\xrightarrow{\text{row-wise lookup}}
    \bm Z,\\
    (\bm\alpha,\bm X)
    &\xrightarrow{\text{source lookup}}
    (\bm\tau,\bm X',\bm\omega').
\end{aligned}
\end{equation}

\mysmallskip\noindent
\textbf{Compiling the scales.}
For the reverse tensor span, a fixed zero bit contributes $1$, while a fixed
one bit contributes the corresponding point coordinate.  The prover
materializes and commits to one prefix matrix
$\bm P^{\mathsf{cmp}}\in\Field^{T\times2m}$.  Its rows start from
$\bm\omega'$ and directly absorb the factors selected by the registered
$\bm D$: for every $j\in[T]$,
\begin{equation}
\label{eq:eta-prefix-recurrence}
\begin{aligned}
P^{\mathsf{cmp}}_{j,1}
&=
\omega'_j
\bigl(1-D_{j,1}+D_{j,1}X'_{j,1}\bigr),\\
P^{\mathsf{cmp}}_{j,h+1}
&=
P^{\mathsf{cmp}}_{j,h}\\[-1mm]
&\quad\cdot
\bigl(1-D_{j,h+1}+D_{j,h+1}X'_{j,h+1}\bigr),
&&h\in[2m-1],\\
\eta_j&=P^{\mathsf{cmp}}_{j,2m}.
\end{aligned}
\end{equation}
For row $j$, $\bm D$ selects exactly the coordinates fixed to one by the
aligned row and column offsets.  Equation~\eqref{eq:alignment-indexed-zero}
ensures that none of the local coordinates contained in $\bm z_j$ is also
treated as fixed.  After the commitments to $\bm X'$, $\bm\omega'$,
$\bm P^{\mathsf{cmp}}$, and $\bm\eta$ are fixed, one transcript challenge
batches the initialization, all prefix steps, and the terminal equality in
Eq.~\eqref{eq:eta-prefix-recurrence}.  Thus, the factors
$\bm1-\bm D+\bm D\circ\bm X'$ are never separately committed, and no
independent row-product proof is required.

Consequently, for every $j\in[T]$,
\begin{equation}
\label{eq:single-bucket-compilation}
    \omega^{\tau_j-1}
    \left.
    (\bm l_{\tau_j}\kronecker\bm r_{\tau_j})
    \right|_j
    =
    \eta_j\tensorspan{\bm z_j}.
\end{equation}
Summing over all split submatrices gives
\begin{equation}
\label{eq:compiled-bva-tbv}
\begin{aligned}
\sum_{\tau=1}^{t}
    \omega^{\tau-1}u_\tau
&=
\sum_{j=1}^{T}
    \eta_j
    \tensorproj{\bm a^{(j)}}{\bm z_j}\\
&= 
\langle\bm a,\bm\zeta\rangle,
\end{aligned}
\end{equation}
where $\bm a$ is the BVA of the split submatrices and
\[
    \bm\zeta
    =
    \bigl(
        N,T,
        \{(k_j,\eta_j,\bm z_j)\}_{j\in[T]}
    \bigr).
\]

The coordinate lookup has $2Tm$ queries, the source-row lookup has $T$
queries, and the batched prefix recurrence involves
$O(Tm)$ entries.  The compilation therefore takes
\[
    O(T\log K)
    =
    O(t\log^3K)
\]
prover time.  Its proof size and verifier time are inherited from the
underlying succinct lookup and matrix-reduction protocols.  For commitment
and MLE opening, all matrices are embedded into their next power-of-two
rectangular covers, with the added entries fixed to zero.

\subsection{Commit to and Hide the Shapes}
\label{sec:shape-hiding}

The matrices $\bm\alpha$, $\bm\gamma$, $\bm\Sigma$, and $\bm D$ depend only
on the original matrix shapes and the deterministic splitting rule.  They are
independent of the evaluation points, so the prover commits during
registration to:
\begin{equation}
\label{eq:shape-commitments}
\begin{aligned}
    C_\alpha&=\pc.\Com(\pp,\bm\alpha),
    &C_\gamma&=\pc.\Com(\pp,\bm\gamma),\\
    C_\Sigma&=\pc.\Com(\pp,\bm\Sigma),
    &C_D&=\pc.\Com(\pp,\bm D).
\end{aligned}
\end{equation}

The one-time registration proof establishes that the committed metadata
describes the deterministic splitting used to construct the BVA.  In
particular, it verifies
\begin{equation}
\label{eq:shape-registration}
\begin{aligned}
    \gamma_{l,j}+\gamma_{r,j}
    &=\log_2k_j,\\
    2^{\gamma_{l,j}}
    &\mid\alpha_{l,j},\\
    2^{\gamma_{r,j}}
    &\mid\alpha_{r,j},
    \qquad j\in[T],
\end{aligned}
\end{equation}
and proves that the resulting rectangles are disjoint and form the
prescribed partition of each source matrix.  The same proof establishes
Eqs.~\eqref{eq:Sigma-gamma-polynomial},
\eqref{eq:D-alpha-consistency}, and
\eqref{eq:alignment-indexed-zero}, thereby binding $C_\Sigma$ and $C_D$ to
the unique matrices derived from $C_\gamma$ and $C_\alpha$.  These checks are
cached with the registration result; subsequent openings reuse $C_\Sigma$
and $C_D$ without repeating the shape-only relations.

Algorithm~\ref{alg:shapeencoding} makes both directions explicit: its
\textsf{Encode} part computes $(\bm\alpha,\bm\gamma)$ from the original
matrix shapes, after which Eqs.~\eqref{eq:sigma-row} and
\eqref{eq:compiled-point-matrices} derive $(\bm\Sigma,\bm D)$.  Its
\textsf{Recover} part reconstructs the original shapes from
$(\bm\alpha,\bm\gamma)$.

% !TEX root = main.tex
\begin{algorithm}[!t]
\caption{Shape Encoding and Shape Recovery}
\label{alg:shapeencoding}

\KwIn{
(\textsf{Encode}) original matrix shapes
$\{(\mu_\tau,\nu_\tau)\}_{\tau\in[t]}$;\\
\qquad the deterministic splitting rule of
Lemma~\ref{lemma:mat-splitting};\\
\qquad (\textsf{Recover}) shape encoding
$\bm\alpha=\concat{\bm\tau,\bm\alpha_l,\bm\alpha_r}\in\Field^{T\times3}$
and
$\bm\gamma=\concat{\bm\gamma_l,\bm\gamma_r}\in\Field^{T\times2}$
}
\KwOut{
(\textsf{Encode}) $T$, $\bm\alpha$, $\bm\gamma$, and
$\{k_j\}_{j\in[T]}$;\\
\qquad (\textsf{Recover}) recovered shapes
$\{(\mu_\tau,\nu_\tau)\}_{\tau\in[t]}$
}

\BlankLine
\textbf{\textsf{Encode}: compute $(\bm\alpha,\bm\gamma)$ from the original shapes.}\\
Initialize $\mathsf{records}\leftarrow[\,]$\;
\For{$\tau\leftarrow1$ \KwTo $t$}{
    Run the deterministic splitting rule of
    Lemma~\ref{lemma:mat-splitting} on the
    $\mu_\tau\times\nu_\tau$ matrix $\bm C_\tau$\;
    \For{each emitted aligned submatrix}{
        Read its row offset, row exponent, column offset, and column exponent as
        $(\alpha_l,\gamma_l,\alpha_r,\gamma_r)$\;
        $k\leftarrow2^{\gamma_l+\gamma_r}$\;
        Append
        $(\tau,\alpha_l,\gamma_l,\alpha_r,\gamma_r,k)$
        to $\mathsf{records}$\;
    }
}
Sort $\mathsf{records}$ by non-increasing $k$, breaking ties by
$(\tau,\alpha_l,\alpha_r)$\;
$T\leftarrow|\mathsf{records}|$\;
\For{$j\leftarrow1$ \KwTo $T$}{
    Read the $j$-th record as
    $(\tau_j,\alpha_{l,j},\gamma_{l,j},
      \alpha_{r,j},\gamma_{r,j},k_j)$\;
}
$\bm\alpha\leftarrow\concat{\bm\tau,\bm\alpha_l,\bm\alpha_r}$; \quad
$\bm\gamma\leftarrow\concat{\bm\gamma_l,\bm\gamma_r}$\;
\Return $(T,\bm\alpha,\bm\gamma,\{k_j\}_{j\in[T]})$\;

\BlankLine
\textbf{\textsf{Recover}: recover the original shapes from $(\bm\alpha,\bm\gamma)$.}\\
$t\leftarrow\max_{j\in[T]}\tau_j$\;
\For{$\tau\leftarrow1$ \KwTo $t$}{
    $\mu_\tau\leftarrow0$; \quad $\nu_\tau\leftarrow0$\;
}
\For{$j\leftarrow1$ \KwTo $T$}{
    $\mu_{\tau_j}\leftarrow
    \max\{\mu_{\tau_j},\alpha_{l,j}+2^{\gamma_{l,j}}\}$\;
    $\nu_{\tau_j}\leftarrow
    \max\{\nu_{\tau_j},\alpha_{r,j}+2^{\gamma_{r,j}}\}$\;
}
\Return $\{(\mu_\tau,\nu_\tau)\}_{\tau\in[t]}$\;
\end{algorithm}


\begin{proposition}[Injectivity of the Shape Encoding]
\label{prop:shape-encoding}
Fix the public BVA layout and the deterministic splitting rule of
Lemma~\ref{lemma:mat-splitting}.  Any well-formed tuple
\[
    (\bm\alpha,\bm\gamma,\bm\Sigma,\bm D)
\]
that satisfies the registration relations uniquely determines the complete
dyadic splitting of every $\bm C_\tau$ and hence its exact shape.
\end{proposition}

\begin{proof}
Consider the row corresponding to the $j$-th BVA bucket.  The first entry of
$\bm\alpha_{j,*}$ identifies its source matrix $\bm C_{\tau_j}$.  The
remaining entries give the row and column offsets
$\alpha_{l,j}$ and $\alpha_{r,j}$, while $\bm\gamma_{j,*}$ gives the
submatrix dimensions
\[
    2^{\gamma_{l,j}}
    \times
    2^{\gamma_{r,j}}.
\]
Therefore, row $j$ uniquely determines the rectangle
\[
\begin{aligned}
&
[\alpha_{l,j},
  \alpha_{l,j}+2^{\gamma_{l,j}})
\\[-1mm]
&\qquad\times
[\alpha_{r,j},
  \alpha_{r,j}+2^{\gamma_{r,j}})
\end{aligned}
\]
inside $\bm C_{\tau_j}$.

For every $\tau\in[t]$, the registration proof establishes that the
rectangles satisfying $\tau_j=\tau$ form the deterministic disjoint
partition of $\bm C_\tau$.  Its exact dimensions are consequently recovered
as
\[
\begin{aligned}
\max_{j:\tau_j=\tau}
\left\{
\alpha_{l,j}+2^{\gamma_{l,j}}
\right\},
\qquad
\max_{j:\tau_j=\tau}
\left\{
\alpha_{r,j}+2^{\gamma_{r,j}}
\right\}.
\end{aligned}
\]
Thus, $(\bm\alpha,\bm\gamma)$ uniquely determines every original matrix
shape and the assignment of every BVA bucket.

Finally, Eq.~\eqref{eq:Sigma-gamma-polynomial} uniquely derives
$\bm\Sigma$ from $\bm\gamma$, while
Eq.~\eqref{eq:D-alpha-consistency} uniquely derives $\bm D$ from the offset
columns of $\bm\alpha$.  The commitments to these matrices therefore bind
the reusable shape-only compiler tables without introducing another choice.
\end{proof}

By PCS binding, the commitments in Eq.~\eqref{eq:shape-commitments}
determine a unique tuple $(\bm\alpha,\bm\gamma,\bm\Sigma,\bm D)$ and hence a unique
collection of matrix shapes by
Proposition~\ref{prop:shape-encoding}.  By PCS hiding and the zero knowledge
of the registration, consistency, lookup, and grand-product proofs, the
verifier learns neither the committed metadata nor the exact matrix shapes
beyond the public BVA layout.

\subsection{The Final TBV Evaluation}
\label{sec:final-tbv-evaluation}

Packed EO-LiteBullet reduces the BVA--TBV inner product to an MLE claim on
each input vector.  For the TBV side, its challenges define the point
\begin{equation}
\label{eq:final-tbv-point}
\bm y^\top
=
\bigl(
y_0,y_0^2,y_0^4,\ldots,y_0^{\kappa/2},
	y_1,\ldots,y_{\log_2(K/\kappa)}
\bigr)
\in\Field^{1\times m}.
\end{equation}
The inverse point $\bm y^{-1}$ is used for the BVA-side opening; the remaining
TBV claim is
\[
    \tensorproj{\bm\zeta}{\bm y}
    =
    \left\langle
        \bm\zeta,
        \tensorspan{\bm y}
    \right\rangle .
\]
We evaluate this claim directly from the compact description of
$\bm\zeta$, without expanding either tensor-structured vector.

\mysmallskip\noindent
\textbf{Global bucket positions.}
For the $j$-th BVA bucket, define
\begin{equation}
\label{eq:global-bucket-gamma}
    \bm\gamma'
    :=
    \bm\gamma_l+\bm\gamma_r,
    \qquad
    \gamma'_j
    =
    \log_2 k_j.
\end{equation}
Thus, exponentiation componentwise gives
\[
    2^{\bm\gamma'}
    =
    (k_1,\ldots,k_T)^\top.
\]
Let
\[
    \bm\alpha'
    =
    (\alpha'_1,\ldots,\alpha'_T)^\top
\]
contain the starting positions of the BVA buckets.  Since the buckets are
concatenated in their fixed BVA order,
\begin{equation}
\label{eq:global-bucket-offsets}
\begin{aligned}
    \alpha'_1&=0,\\
    \alpha'_{j+1}
        &=\alpha'_j+2^{\gamma'_j},
        &&j\in[T-1],\\
    \alpha'_T+2^{\gamma'_T}
        &=N.
\end{aligned}
\end{equation}

Both $\bm\gamma'$ and $\bm\alpha'$ are determined by the public ordered
bucket sizes in $\mathcal L$.  The verifier therefore derives them directly
from Eqs.~\eqref{eq:global-bucket-gamma} and
\eqref{eq:global-bucket-offsets}; neither vector is committed or opened.

\mysmallskip\noindent
\textbf{Compiling the restriction of $\tensorspan{\bm y}$.}
From the public $\bm\gamma'$ and $\bm\alpha'$, the verifier derives
\[
    \bm\Sigma',\bm D'
    \in\Field^{T\times m},
\]
while the prover constructs and commits to
$\bm Z'\in\Field^{T\times m}$.  For every $j\in[T]$, the $j$-th row of
$\bm\Sigma'$ is
\begin{equation}
\label{eq:Sigma-prime-row}
    \bm\Sigma'_{j,*}
    =
    \bigl(
        1,\ldots,\gamma'_j,
        0,\ldots,0
    \bigr).
\end{equation}
Thus, its nonzero entries identify the first $\gamma'_j$ coordinates of
$\bm y$.  The $j$-th row of $\bm D'$ is the $m$-bit little-endian
decomposition of $\alpha'_j$, while
\begin{equation}
\label{eq:Z-prime-row}
    \bm Z'_{j,*}
    =
    \bigl(
        y_1,\ldots,y_{\gamma'_j},
        0,\ldots,0
    \bigr).
\end{equation}

Since the ordered bucket sizes and offsets are public, these definitions
uniquely determine $\bm\Sigma'$ and $\bm D'$.  Bucket alignment also makes
the first $\gamma'_j$ bits of $\alpha'_j$ zero by construction.  Thus, no
commitment, Booleanity proof, reconstruction proof, or alignment lookup is
needed for either public table.

Finally, $\bm Z'$ is selected from $\bm y$ according to $\bm\Sigma'$.
After committing to $\bm Z'$, the prover gives the genuine row-tagged LogUp:
\begin{equation}
\label{eq:Z-prime-lookup}
\begin{aligned}
\left\{
(j,Z'_{j,h},\Sigma'_{j,h})
:
\substack{j\in[T],\\h\in[m]}
\right\}
&\quad\subseteq
\left\{(j,0,0):j\in[T]\right\}\\
&\qquad\cup
\left\{
(j,y_h,h)
:
\substack{j\in[T],\\h\in[m]}
\right\}.
\end{aligned}
\end{equation}
The public row tag prevents a query from borrowing an entry from another
bucket.  Since $\bm\Sigma'$ is verifier-derived, this lookup proves that row
$j$ of committed $\bm Z'$ consists of the first $\gamma'_j$ coordinates of
$\bm y$, followed by zeros.

\mysmallskip\noindent
\textbf{Compiling the tensor scales.}
The bit decomposition $\bm D'$ records the coordinates fixed by the global
offset of each bucket.  A fixed zero bit contributes $1$ to the reverse
tensor span, while a fixed one bit contributes the corresponding coordinate
of $\bm y$.  Since both $\bm D'$ and $\bm y$ are public, the verifier directly
derives:
\begin{equation}
\label{eq:eta-prime-public}
    \bm\eta'
    =
    \mathsf{GrandProdRow}\!\left(
        \bm1_{T\times m}
        -
        \bm D'
        +
        \bm D'
        \circ
        (\bm1_T\bm y^\top)
    \right).
\end{equation}
No commitment or grand-product proof is needed for $\bm\eta'$.  Bucket
alignment ensures that the coordinates contained in $\bm Z'$ are not also
included in the scale.  Consequently, for every bucket $j$:
\begin{equation}
\label{eq:y-bucket-restriction}
    \left.
        \tensorspan{\bm y}
    \right|_j
    =
    \eta'_j
    \tensorspan{
        \bm Z'_{j,[1:\gamma'_j]}
    }.
\end{equation}

\mysmallskip\noindent
\textbf{Final evaluation.}
Recall that the $j$-th bucket of $\bm\zeta$ is
\[
    \eta_j\tensorspan{\bm z_j}.
\]
Since $k_j\le N$, we have $\gamma'_j\le m$, and all genuine entries of
$\bm z_j$ lie in the first $m$ columns of $\bm Z$.  Writing
\[
\bm\eta=(\eta_1,\ldots,\eta_T)^\top,
\]
the mixed-product property and
Eq.~\eqref{eq:y-bucket-restriction} give
\begin{equation}
\label{eq:final-tbv-evaluation}
\begin{aligned}
\tensorproj{\bm\zeta}{\bm y}
&=
\sum_{j=1}^{T}
\eta_j\eta'_j
\prod_{h=1}^{\gamma'_j}
\left(
    1+Z_{j,h}Z'_{j,h}
\right).
\end{aligned}
\end{equation}
The zero padding in $\bm Z$ and $\bm Z'$ contributes only unit factors, so
the row-wise product is exactly over the
$\gamma'_j=\log_2k_j$ local coordinates of each bucket.

The proof does not materialize separate factor and scale commitments and does
not invoke a second row-product protocol.  Instead, it commits to a prefix
matrix $\bm P\in\Field^{T\times m}$ and a bucket-evaluation vector
$\bm v\in\Field^T$ satisfying the single batched recurrence:
\begin{equation}
\label{eq:final-prefix-recurrence}
\begin{aligned}
P_{j,1}
    &=\eta_j\eta'_j,\\
P_{j,h+1}
    &=P_{j,h}\left(1+Z_{j,h}Z'_{j,h}\right),
    &&j\in[T],\ h\in[m-1],\\
v_j
    &=P_{j,m}\left(1+Z_{j,m}Z'_{j,m}\right),
    &&j\in[T],\\
\tensorproj{\bm\zeta}{\bm y}
    &=\bm1_T^\top\bm v.
\end{aligned}
\end{equation}
One transcript challenge batches the initialization, all prefix steps, and
the final output into one degree-three hidden relation over the existing
commitments to $\bm Z$ and $\bm\eta$, the LogUp-bound commitment to
$\bm Z'$, the public $\bm\eta'$, and the new commitments $C_P$ and $C_v$.
The prover then opens $C_v$ at the all-half point of its public padded row
cover.  Multiplying the hidden opening by the row-cover size yields
$\bm1_T^\top\bm v$, which is checked against the hidden terminal scalar from
Packed EO-LiteBullet.

All matrices introduced in this subsection have dimensions at most
$T\times m$.  Deriving the public tables, constructing the prefixes, proving
the $\bm Z'$ LogUp, and proving the recurrence therefore take
\[
    O(Tm)
    =
    O(T\log K)
    =
    O(t\log^3K)
\]
prover time.  The final proof adds three primary commitments
$(C_{Z'},C_P,C_v)$; $\bm\Sigma'$ and $\bm\eta'$ require neither a commitment
nor a PCS opening.  The proof size and verifier time are inherited from the
underlying succinct lookup and matrix-reduction protocols.

\subsection{The Complete Protocol and Complexity}
\label{sec:complete-protocol}

The protocol consists of a one-time registration phase and an opening phase
that is repeated for each collection of matrix-evaluation claims.  The
registration phase commits to the matrices and their hidden shapes.  During
an opening, the original matrix points are compiled into one TBV, after which
Packed EO-LiteBullet reduces the resulting BVA--TBV inner product to one
opening of the committed BVA and one compact TBV evaluation.

Algorithm~\ref{alg:mleclaimconversion} summarizes this complete MLE claim
conversion argument.

% !TEX root = main.tex
%\colorbox{gray!5!white}{%
\begin{algorithm}[!t]
\caption{MLE Claim Conversion Argument}
\label{alg:mleclaimconversion}
\small

\KwIn{
A base PCS $\pc$ over $\Field$ with public parameters $\pp$;
packed-fold factor $\kappa=2^\tau$\;
\qquad (\textsf{Register}) matrices
$\{\bm C_\tau\}_{\tau\in[t]}$\;
\qquad (\textsf{Open}) public claimed values $\{u_\tau\}_{\tau\in[t]}$
and prover-private evaluation points
$\{\bm x_{l,\tau},\bm x_{r,\tau}\}_{\tau\in[t]}$
}
\KwOut{
\qquad (\textsf{Collection object})
$(C_a,C_\alpha,C_\gamma,C_\Sigma,C_D,\pi_{\mathsf{reg}})$\;
\qquad (\textsf{Opening proof}) $(C_X,\pi)$\;
\qquad (verification) $\TRUE/\FALSE$
}

\BlankLine
\textbf{Register.}\\
Apply the deterministic splitting rule of
Lemma~\ref{lemma:mat-splitting} to $\{\bm C_\tau\}_{\tau\in[t]}$\;
Obtain the BVA buckets
$\{\bm a^{(j)}\}_{j\in[T]}$, bucket sizes $\{k_j\}_{j\in[T]}$,
the logical length $K$, and the BVA $\bm a$ in
Eq.~\eqref{eq:matrix-bva}\;
Run \textsf{Encode} of Algorithm~\ref{alg:shapeencoding} on
$\{(\mu_\tau,\nu_\tau)\}_{\tau\in[t]}$ to obtain
$\bm\alpha$ and $\bm\gamma$, and derive the shape-only matrices
$\bm\Sigma$ and $\bm D$\;
$C_\alpha\leftarrow\pc.\Com(\pp,\bm\alpha)$; \quad
$C_\gamma\leftarrow\pc.\Com(\pp,\bm\gamma)$; \quad
$C_\Sigma\leftarrow\pc.\Com(\pp,\bm\Sigma)$; \quad
$C_D\leftarrow\pc.\Com(\pp,\bm D)$; \quad
$C_a\leftarrow\pc.\Com(\pp,\bm a)$\;
Construct the one-time registration proof $\pi_{\mathsf{reg}}$ for
Eqs.~\eqref{eq:shape-registration}--\eqref{eq:alignment-indexed-zero} and
the partition consistency of the split submatrices\;
\Return $(C_a,C_\alpha,C_\gamma,C_\Sigma,C_D,\pi_{\mathsf{reg}})$\;


\BlankLine
\textbf{Open.}\\
Form the point matrix $\bm X$ in Eq.~\eqref{eq:combined-point-matrix} and
send only $C_X\leftarrow\pc.\Com(\pp,\bm X)$\;
$\VtoP$: sample $\omega \gets \Field^*$\;
$c\leftarrow\sum_{\tau=1}^{t}\omega^{\tau-1}u_\tau$\;
Construct the auxiliary matrices
$\bm X',\bm Z,\bm\omega',\bm P^{\mathsf{cmp}},\bm\eta$ using the registered
$\bm\Sigma,\bm D$, commit to these opening-time witnesses, and form the compact TBV
$\bm\zeta=(N,T,\{(k_j,\eta_j,\bm z_j)\}_{j\in[T]})$\;
Construct the point-compilation proof $\pi_{\mathsf{cmp}}$ for
Eq.~\eqref{eq:complete-point-compilation} using the cached alignment check and
one batched prefix recurrence from $\bm\omega'$ to $\bm\eta$\;
$\pi_{\mathsf{ip}}\leftarrow
\mathsf{PackedEO\text{-}LiteBullet}.\Prove(K,\kappa,c,\bm a,\bm\zeta)$\;
Parse $\pi_{\mathsf{ip}}$ to obtain $\bm y$ and the hidden terminal claims
$\bar a=\tensorproj{\bm a}{\bm y^{-1}}$ and
$\bar b=\tensorproj{\bm\zeta}{\bm y}$\;
$\pi_a\leftarrow\pc.\Open(\pp,C_a,\bar a,\bm y^{-1},\bm a)$\;
Derive the public tables
$\bm\gamma',\bm\alpha',\bm\Sigma',\bm D',\bm\eta'$ and commit only to
the hidden $\bm Z'$, prefix matrix $\bm P$, and bucket-evaluation vector
$\bm v$ as in Section~\ref{sec:final-tbv-evaluation}\;
Construct the final TBV evaluation proof $\pi_{\mathsf{tbv}}$ for
Eq.~\eqref{eq:final-prefix-recurrence}, including the row-tagged LogUp for
$\bm Z'$\;
$\pi\leftarrow(\pi_{\mathsf{cmp}},\pi_{\mathsf{ip}},\pi_a,\pi_{\mathsf{tbv}})$\;
\Return $(C_X,\pi)$\;

\BlankLine
\textbf{Verify.}\\
Check $\pi_{\mathsf{reg}}$ once and cache the result\;
Parse $\pi$ as
$(\pi_{\mathsf{cmp}},\pi_{\mathsf{ip}},\pi_a,\pi_{\mathsf{tbv}})$\;
Verify $\pi_{\mathsf{cmp}}$ against
$C_\alpha$, $C_\gamma$, $C_\Sigma$, $C_D$, and $C_X$, without receiving
the entries of $\bm X$\;
Verify $\pi_{\mathsf{ip}}$ for the claim
$c=\langle\bm a,\bm\zeta\rangle$ and derive $\bm y$ together with the
commitments to $\bar a$ and $\bar b$\;
Verify the BVA opening
$\pc.\Verify(\pp,C_a,\bar a,\bm y^{-1},\pi_a)$\;
Verify $\pi_{\mathsf{tbv}}$ for the claim
$\bar b=\tensorproj{\bm\zeta}{\bm y}$\;
\Return $\TRUE$ iff all checks accept\;
\end{algorithm}


\mysmallskip\noindent
\textbf{Registration.}
The prover applies the deterministic splitting rule of
Lemma~\ref{lemma:mat-splitting} to
$\{\bm C_\tau\}_{\tau\in[t]}$, obtaining the $T$ power-of-two submatrices.
Their row-major vectorizations are concatenated in the prescribed order to
form the BVA $\bm a$.  For every BVA bucket, the prover also constructs
\[
    \bm\alpha
    =
    \concat{\bm\tau,\bm\alpha_l,\bm\alpha_r},
    \qquad
    \bm\gamma
    =
    \concat{\bm\gamma_l,\bm\gamma_r}.
\]
It also derives the shape-only matrices $\bm\Sigma$ and $\bm D$ from
Eqs.~\eqref{eq:sigma-row} and
\eqref{eq:compiled-point-matrices}.  It publishes the shape commitments
before the dense BVA commitment:
\[
\begin{aligned}
    C_\alpha&=\pc.\Com(\pp,\bm\alpha),
    &C_\gamma&=\pc.\Com(\pp,\bm\gamma),\\
    C_\Sigma&=\pc.\Com(\pp,\bm\Sigma),
    &C_D&=\pc.\Com(\pp,\bm D),\\
    C_a&=\pc.\Com(\pp,\bm a).
\end{aligned}
\]

The one-time registration proof establishes that these commitments encode the
same deterministic split.  In particular, it proves that the committed BVA
is the concatenation of the split submatrices and that, for every
$j\in[T]$,
\[
\begin{aligned}
    \gamma_{l,j}+\gamma_{r,j}
        &=\log_2k_j,\\
    2^{\gamma_{l,j}}
        &\mid\alpha_{l,j},\\
    2^{\gamma_{r,j}}
        &\mid\alpha_{r,j}.
\end{aligned}
\]
It additionally proves that the rectangles encoded by
$(\bm\alpha,\bm\gamma)$ form the prescribed disjoint partition of each
source matrix.  The same registration proof binds $\bm\Sigma$ to
$\bm\gamma$, binds $\bm D$ to $\bm\alpha$, and proves their alignment through
Eqs.~\eqref{eq:Sigma-gamma-polynomial}--\eqref{eq:alignment-indexed-zero}.
This proof is verified once, and all five collection commitments are reused
for subsequent openings.

\mysmallskip\noindent
\textbf{Point compilation.}
For one opening batch, the claimed values
$\{u_\tau\}_{\tau\in[t]}$ and the corresponding points
$\{\bm x_{l,\tau},\bm x_{r,\tau}\}_{\tau\in[t]}$ are fixed by the prover
before the batching challenge.  The prover forms the point matrix $\bm X$ in
Eq.~\eqref{eq:combined-point-matrix} and sends only its hiding commitment:
\[
    C_X=\pc.\Com(\pp,\bm X).
\]
The verifier does not receive the entries of $\bm X$ or the original
evaluation points.  It samples $\omega\ChooseRandom\Field^*$ after $C_X$ is
fixed.

Using the registered $\bm\Sigma$ and $\bm D$, the committed primary metadata,
and the challenge $\omega$, the prover
constructs
\[
    \bm X',\quad
    \bm Z,\quad
    \bm\omega',\quad
    \bm P^{\mathsf{cmp}},\quad
    \bm\eta.
\]
Their consistency is established by
\begin{equation}
\label{eq:complete-point-compilation}
\begin{aligned}
    (\bm X',\bm\Sigma)
    &\xrightarrow{\text{Eq.~\eqref{eq:coordinate-lookup}}}
        \bm Z,\\
    (\bm\alpha,\bm X)
    &\xrightarrow{\text{Eq.~\eqref{eq:source-row-lookup}}}
        (\bm\tau,\bm X',\bm\omega'),\\
    (\bm\omega',\bm D,\bm X')
    &\xrightarrow{\text{Eq.~\eqref{eq:eta-prefix-recurrence}}}
        (\bm P^{\mathsf{cmp}},\bm\eta).
\end{aligned}
\end{equation}
The cached indexed-zero relation in
Eq.~\eqref{eq:alignment-indexed-zero} already proves that the local
coordinates selected by registered $\bm\Sigma$ are disjoint from the fixed
coordinates selected by registered $\bm D$.  All opening-time relations are
zero knowledge, so the compilation binds to $C_X$ without disclosing its
entries.

These relations compile the original claims into
\begin{equation}
\label{eq:complete-bva-tbv}
\begin{aligned}
    \sum_{\tau=1}^{t}
        \omega^{\tau-1}u_\tau
    &=
    \sum_{j=1}^{T}
        \eta_j
        \tensorproj{\bm a^{(j)}}{\bm z_j}\\
    &=
    \langle\bm a,\bm\zeta\rangle,
\end{aligned}
\end{equation}
where
\[
    \bm\zeta
    =
    \bigl(
        N,T,
        \{(k_j,\eta_j,\bm z_j)\}_{j\in[T]}
    \bigr).
\]

\mysmallskip\noindent
\textbf{BVA--TBV reduction.}
The prover runs Packed EO-LiteBullet on
Eq.~\eqref{eq:complete-bva-tbv}.  Its transcript reduces the aggregated
claim to
\begin{equation}
\label{eq:complete-packed-reduction}
    \tensorproj{\bm a}{\bm y^{-1}},
    \qquad
    \tensorproj{\bm\zeta}{\bm y}.
\end{equation}
The first claim is an MLE opening of the original BVA.  The prover therefore
opens $C_a$ at $\bm y^{-1}$ using the underlying PCS, without materializing
the logical padding of $\bm a$.

To prove the second claim, the prover constructs
\[
    \bm Z',\quad
    \bm P,\quad
    \bm v,
\]
as in Section~\ref{sec:final-tbv-evaluation}.  The verifier derives
$\bm\gamma'$, $\bm\alpha'$, $\bm\Sigma'$, $\bm D'$, and $\bm\eta'$ directly
from the public layout and $\bm y$.  The row-tagged LogUp in
Eq.~\eqref{eq:Z-prime-lookup} binds committed $\bm Z'$ to public
$\bm\Sigma'$, while the single relation in
Eq.~\eqref{eq:final-prefix-recurrence} binds $C_P$ and $C_v$ to the existing
commitments to $\bm Z$ and $\bm\eta$.  Opening $C_v$ at the all-half point
and checking its hidden sum against the Packed EO terminal scalar completes
the claim.  Thus, the verifier receives neither the expanded length-$N$
representation of $\bm\zeta$ nor the original evaluation points.

Every opening-time witness matrix is committed before the challenge used in
its lookup or polynomial check.  Public-derived tables enter those checks
directly and require no PCS opening.  The verifier accepts if the cached
registration proof, the point-compilation proof, the Packed EO-LiteBullet
transcript, the separate BVA opening, and the final TBV evaluation proof all
accept.

\mysmallskip\noindent
\textbf{Prover time.}
Recall that
\[
    m=\log_2K,
    \qquad
    T\le t(m+1)^2.
\]
The BVA contains exactly $N$ matrix entries.  Constructing it and computing
its commitment take $O(N)$ time.  Constructing the shape metadata and proving
its consistency require $O(Tm)$ additional work.  Hence, registration takes
\[
    O(N+Tm)
    =
    O(N+t\log^3K)
\]
prover time.

For each opening, constructing $\bm X$ takes $O(tm)$ time.  All matrices in
the point-compilation proof have $O(Tm)$ entries, so
Eq.~\eqref{eq:complete-point-compilation} requires $O(Tm)$ prover time, while
the shape-only checks for $\bm\Sigma$ and $\bm D$ are not repeated.  The
public final tables and the prefix recurrence also have dimensions bounded by
$T\times m$ and require another $O(Tm)$ operations.

Packed EO-LiteBullet takes
\[
    O\left(
        \frac{N}{\kappa}M(\kappa)
        +
        \frac{N}{\kappa}
        +
        T
    \right)
\]
time, where $M(\kappa)$ is the cost of multiplying degree-$(\kappa-1)$
polynomials.  The PCS opening of the BVA scans its $N$ stored entries.
Therefore, the total opening-prover time is
\begin{equation}
\label{eq:complete-prover-time}
    O\left(
        N
        +
        \frac{N}{\kappa}M(\kappa)
        +
        T\log K
    \right).
\end{equation}
Using $T\le t(\log_2K+1)^2$, this becomes
\[
    O\left(
        N
        +
        \frac{N}{\kappa}M(\kappa)
        +
        t\log^3K
    \right).
\]
For fixed $\kappa$, the opening-prover time is
\[
    O(N+t\log^3K).
\]

\mysmallskip\noindent
\textbf{Memory.}
The persistent representation stores the $N$ original matrix entries and the
$O(Tm)$ registered shape tables, including $\bm\Sigma$ and $\bm D$.  The
opening-time point-compilation and final-prefix matrices also contain
$O(Tm)$ field elements.  Hence, for field-valued inputs, the total storage is
\[
    O(N+T\log K).
\]
When the packed first round streams native short-bitwidth entries, the
full-field working memory after this round is
\[
    O\left(
        \frac{N}{\kappa}
        +
        \kappa
        +
        T\log K
    \right),
\]
in addition to the native representation of the matrices.  No expanded TBV
or padding proportional to a larger logical cover is materialized.

\mysmallskip\noindent
\textbf{Proof size and verifier time.}
Packed EO-LiteBullet contributes
\[
    2\kappa-2
    +
    2\log_2(K/\kappa)
    +
    2
\]
field elements.  The remaining consistency, lookup, grand-product,
matrix-reduction, and PCS-opening proofs have logarithmic communication and
verification cost in their matrix dimensions.  Since
\[
    T\le t(\log_2K+1)^2
    \qquad\text{and}\qquad
    t\le N,
\]
the total per-opening proof size and cryptographic verifier time are
\[
    O(\kappa+\log K).
\]
For fixed $\kappa$, both are $O(\log K)$.  The registration proof has the
same asymptotic size and is verified only once.

If the values $\{u_\tau\}_{\tau\in[t]}$ are supplied explicitly, reading and
aggregating them adds $O(t)$ verifier time.  When these values are already
committed, their weighted sum under $\omega$ can instead be verified by one
MLE opening, preserving the $O(\kappa+\log K)$ verifier bound.

For a non-power-of-two total size, the protocol replaces $N$ by its
power-of-two BVA cover.  Since this cover is smaller than $2N$, all preceding
asymptotic bounds remain unchanged.





